\documentclass[a4paper,11pt]{article}
\usepackage[utf8]{inputenc}
\usepackage[spanish]{babel}
\usepackage{geometry}
\usepackage{enumitem}
\usepackage{sectsty}
\usepackage{titlesec}
\usepackage{tikz}
\usepackage{array}

\geometry{
    a4paper,
    left=20mm,
    top=18mm,
    right=20mm,
    bottom=20mm
}

% Sin números de página
\pagestyle{empty}

% Estilo de secciones con números en círculos
\newcommand{\circled}[1]{%
    \tikz[baseline=(char.base)]{%
        \node[shape=circle,draw,inner sep=0.5pt,fill=black,text=white,font=\bfseries\footnotesize] (char) {#1};%
    }%
}

% Formato de secciones
\titleformat{\section}
{\normalfont\Large\bfseries}
{\circled{\thesection}}{0.5em}{}

\titlespacing*{\section}{0pt}{12pt}{8pt}

\titleformat{\subsection}
{\normalfont\normalsize\bfseries}
{}{0em}{}

\titlespacing*{\subsection}{8pt}{4pt}{4pt}

\begin{document}

% Encabezado del curso
\begin{center}
\small{Universidad de Granada}

\vspace{0.3cm}

\textbf{\large{Producción e Interacción Oral - Nivel 7}}

\vspace{0.4cm}

\small
\begin{tabular}{@{}ll@{}}
\textbf{Grupos:} & A (Aula 24) y B (Aula 08) \\
\textbf{Centro:} & Centro de Lenguas Modernas / Edificio A \\
\textbf{Profesor:} & Javier Benítez Lainez \\
\textbf{Tutorías:} & Martes 9:00-10:00 y 14:30-15:30 \\
& Sala de profesores / Email: benitezl@go.ugr.es
\end{tabular}
\end{center}

\vspace{0.5cm}
\section*{Gestión de Preguntas en Q&A}

\vspace{0.4cm}

\section{Guía para Manejar el Periodo de Preguntas y Respuestas}

\vspace{0.2cm}

---

\vspace{0.2cm}

\section{FUNDAMENTOS DEL Q&A}

\vspace{0.2cm}

\subsection{¿Qué es el Q&A?}

El periodo de preguntas y respuestas (Q&A = Questions & Answers) es una parte crítica de cualquier presentación donde interactúas directamente con el audiencia, respondiendo dudas, aclarando puntos y defendiendo tu posición.

\vspace{0.2cm}

\subsection{Objetivos del Q&A}

1. \textbf{Aclarar} - Resolver confusiones o malentendidos

2. \textbf{Profundizar} - Expandir en puntos de interés

3. \textbf{Reforzar} - Reiterar mensajes clave

4. \textbf{Persuadir} - Convencer a escépticos

5. \textbf{Conectar} - Construir relación con el audiencia

\vspace{0.2cm}

---

\vspace{0.2cm}

\section{PREPARACIÓN ANTES DEL Q&A}

\vspace{0.2cm}

\subsection{Anticipación de preguntas}

\vspace{0.2cm}

\section{Tipos de preguntas a anticipar}

\begin{itemize}

  \item \textbf{Preguntas fácticas} - Datos, cifras, evidencia

  \item \textbf{Preguntas de clarificación} - Explicación de conceptos

  \item \textbf{Preguntas críticas} - Desafíos a tu argumento

  \item \textbf{Preguntas "¿y si?"} - Escenarios hipotéticos

  \item \textbf{Preguntas personales} - Tu opinión o experiencia

\end{itemize}

\vspace{0.2cm}

\section{Técnica de "lluvia inversa"}

\begin{itemize}

  \item Haz una lluvia de ideas de las 10 preguntas más probables

  \item Prepara respuestas concisas (30-60 segundos cada una)

  \item Ten datos de apoyo listos para citar

  \item Prepara ejemplos adicionales si te piden elaborar

\end{itemize}

\vspace{0.2cm}

\subsection{Preparación de respuestas}

\vspace{0.2cm}

\section{Estructura de respuesta efectiva}

\begin{itemize}

  \item \textbf{ACKNOWLEDGE} - Reconoce la pregunta

  \item \textbf{ANSWER} - Responde directamente

  \item \textbf{ELABORATE} - Explica brevemente

  \item \textbf{EXAMPLE} - Da un ejemplo si es necesario

  \item \textbf{CHECK} - Verifica si respondiste adecuadamente

\end{itemize}

\vspace{0.2cm}

---

\vspace{0.2cm}

\section{DURANTE EL Q&A}

\vspace{0.2cm}

\subsection{Recibir la pregunta}

\vspace{0.2cm}

\section{Escucha activa}

\begin{itemize}

  \item \textbf{Mantén contacto visual} con el preguntón

  \item \textbf{Asiente} mientras escuchas

  \item \textbf{Toma notas} si la pregunta es compleja

  \item \textbf{No interrumpas} - deja que terminen

\end{itemize}

\vspace{0.2cm}

\section{Repetir o parafrasear la pregunta}

\begin{itemize}

  \item \textbf{"Si les entiendo bien, están preguntando sobre..."}

  \item \textbf{"Déjenme asegurarme de entender correctamente..."}

  \item \textbf{"En otras palabras, quieren saber..."}

\end{itemize}

\vspace{0.2cm}

\textbf{¿Por qué parafrasear?}

\begin{itemize}

  \item Confirms que entendiste

  \item Da tiempo para pensar

  \item Permite a otros escuchar la pregunta

  \item Puedes reencuadrar la pregunta

\end{itemize}

\vspace{0.2cm}

---

\vspace{0.2cm}

\subsection{Clasificación de preguntas}

\vspace{0.2cm}

\section{Preguntas constructivas}

\begin{itemize}

  \item Son legítimas y relevantes

  \item Muestran interés genuino

  \item Piden clarificación o profundización

\end{itemize}

\vspace{0.2cm}

\section{Preguntas críticas}

\begin{itemize}

  \item Desafían tu posición

  \item Piden evidencia adicional

  \item Buscan debilidades

\end{itemize}

\vspace{0.2cm}

\section{Preguntas problemáticas}

\begin{itemize}

  \item No relacionadas con el tema

  \item Muy largas o confusas

  \item Agresivas o con agenda oculta

  \item Intentan dominar o desviar

\end{itemize}

\vspace{0.2cm}

---

\vspace{0.2cm}

\subsection{Responder a diferentes tipos de preguntas}

\vspace{0.2cm}

\section{Preguntas Fáciles/Directas}

\begin{itemize}

  \item Responde directamente y concisamente

  \item "Esa es una excelente pregunta. La respuesta es..."

  \item "Sí, absolutamente. Permítanme explicar..."

  \item "Gracias por preguntar. Básicamente..."

\end{itemize}

\vspace{0.2cm}

\section{Preguntas Complejas/Técnicas}

\begin{itemize}

  \item Divídelas en partes

  \item "Esa pregunta tiene varios aspectos. Primero..., segundo..."

  \item "Buena pregunta. Déjenme abordar esto paso a paso..."

  \item Usa analogías o ejemplos si ayuda

\end{itemize}

\vspace{0.2cm}

\section{Preguntas Críticas/Desafiantes}

\begin{itemize}

  \item Mantén la calma y el respeto

  \item Reconoce la validez del desafío

  \item "Veo su punto. Es una preocupación legítima..."

  \item "Gracias por plantear eso. Permítanme explicar mi posición..."

  \item Usa evidencia y lógica, no emoción

  \item "La evidencia sugiere que..."

  \item "Según los datos que presenté..."

\end{itemize}

\vspace{0.2cm}

\section{Preguntas Hipotéticas ("¿Qué pasa si...?")}

\begin{itemize}

  \item Reconoce la naturaleza hipotética

  \item "En ese escenario hipotético, probablemente..."

  \item "Es una pregunta interesante. Si asumimos que..."

  \item Usa condicionales

  \item "Si X ocurriera, entonces Y..."

  \item No especules excesivamente

  \item "Es difícil predecir con certeza, pero..."

\end{itemize}

\vspace{0.2cm}

\section{Preguntas para las que NO tienes respuesta}

\begin{itemize}

  \item Sé honesto/a

  \item "Esa es una excelente pregunta para la que no tengo la respuesta ahora."

  \item "No tengo ese dato específico, pero puedo decirles que..."

  \item Ofrece seguir

  \item "Permítanme investigar y les daré una respuesta..."

  \item "¿Podemos conversar sobre esto después de la presentación?"

\end{itemize}

\vspace{0.2cm}

\section{Preguntas fuera de tema}

\begin{itemize}

  \item Redirige cortésmente

  \item "Esa es una pregunta interesante, pero no directamente relacionada con mi tema actual."

  \item "Prefiero concentrarnos en el tema de hoy, pero podemos discutir eso después."

  \item Conecta si es posible

  \item "Aunque no es directamente el tema, se relaciona con..."

\end{itemize}

\vspace{0.2cm}

\section{Preguntas hostiles/agresivas}

\begin{itemize}

  \item Mantén la compostura

  \item No respondas con hostilidad

  \item Valida la emoción, no el tono

  \item "Veo que sientes fuertemente sobre esto..."

  \item Reencuadra la pregunta

  \item "Entiendo que hay preocupación. Déjenme explicar..."

  \item No te defiendas excesivamente

  \item Usa datos objetivos

  \item Si continúa la agresión

  \item "Agradezco su pregunta, pero prefiero que mantengamos un tono respetuoso."

  \item "Les invito a conversar sobre esto individualmente después."

\end{itemize}

\vspace{0.2cm}

---

\vspace{0.2cm}

\section{TÉCNICAS DE GESTIÓN}

\vspace{0.2cm}

\subsection{Técnica del "Puente"}

\vspace{0.2cm}

\section{Usar la pregunta como puente a tu mensaje}

\textbf{Pregunta:} "¿Y si su plan no funciona?"

\vspace{0.2cm}

\textbf{Respuesta:} "Esa es una preocupación válida sobre los riesgos \textbf{[PUENTE]} y precisamente por eso hemos incluido múltiples salvaguardas en el plan para asegurar su éxito..."

\vspace{0.2cm}

\textbf{Fórmulas de puente:}

\begin{itemize}

  \item \textbf{"Lo que esto realmente nos lleva a pensar es..."}

  \item \textbf{"Esa pregunta se relaciona con un punto importante..."}

  \item \textbf{"Lo que está detrás de su pregunta es..."}

  \item \textbf{"En realidad, esto nos lleva a considerar..."}

\end{itemize}

\vspace{0.2cm}

\subsection{Técnica de "Reencuadre"}

\vspace{0.2cm}

\section{Cambiar el marco de la pregunta}

\textbf{Pregunta:} "¿No es esto demasiado costoso?"

\vspace{0.2cm}

\textbf{Respuesta:} "Más que costo, lo que estamos considerando es inversión \textbf{[REENCUADRE]} Una inversión que generará retornos significativos a largo plazo..."

\vspace{0.2cm}

\subsection{Técnica de "Posposición"}

\vspace{0.2cm}

\section{Para preguntas muy complejas o fuera de tema}

\begin{itemize}

  \item \textbf{"Esa es una pregunta compleja que merece más tiempo del que tenemos disponible."}

  \item \textbf{"¿Podemos tomar esto offline después de la sesión?"}

  \item \textbf{"Prefiero no especular sin datos concretos. Puedo investigar y responderles luego."}

\end{itemize}

\vspace{0.2cm}

\subsection{Técnica de "Devolución"}

\vspace{0.2cm}

\section{Cuando apropiado, devolver la pregunta}

\begin{itemize}

  \item \textbf{"Eso es algo interesante. ¿Alguien más tiene pensamientos sobre esto?"}

  \item \textbf{"¿Cómo abordarían esto ustedes?"}

  \item \textbf{"Me gustaría escuchar su perspectiva sobre esto."}

\end{itemize}

\vspace{0.2cm}

---

\vspace{0.2cm}

\section{CLAUSURA DEL Q&A}

\vspace{0.2cm}

\subsection{Señales de cierre}

\begin{itemize}

  \item \textbf{"Tenemos tiempo para una pregunta más..."}

  \item \textbf{"¿Alguna última pregunta antes de concluir?"}

  \item \textbf{"Una última intervención..."}

\end{itemize}

\vspace{0.2cm}

\subsection{Si no hay más preguntas}

\begin{itemize}

  \item Agradece al audiencia

  \item "Gracias por sus excelentes preguntas."

  \item "Agradezco su atención y participación."

  \item Refuerza el mensaje principal

  \item "Para cerrar, quiero reiterar que..."

  \item "Llevándonos todo a una conclusión:..."

  \item Proporciona siguiente paso

  \item "Para los interesados, estoy disponible para..."

  \item "Pueden encontrarme en..."

\end{itemize}

\vspace{0.2cm}

\subsection{Si el tiempo se acabó y quedaron preguntas}

\begin{itemize}

  \item Agradece y reconoce

  \item "Lamento no poder responder a todas, pero agradezco su interés."

  \item Ofrece seguimiento

  \item "Los que quedaron con preguntas, puedo responderles individualmente después."

  \item "Pueden enviarme sus preguntas por email."

\end{itemize}

\vspace{0.2cm}

---

\vspace{0.2cm}

\section{ERRORES COMUNES A EVITAR}

\vspace{0.2cm}

\subsection{❌ Errores fatales}

1. \textbf{Defensiveness} - Ponerse a la defensiva

2. \textbf{Evasiveness} - Evadir respuestas

3. \textbf{Hostility} - Responder con hostilidad

4. \textbf{Speculation} - Especular sin datos

5. \textbf{Monopolization} - Dejar que uno domine

6. \textbf{Embarrassment} - Hacer sentir mal al preguntón

7. \textbf{Lying} - Inventar respuestas

8. \textbf{Rambling} - Responder de forma interminable

\vspace{0.2cm}

\subsection{✅ Prácticas excelentes}

1. \textbf{Preparation} - Anticipar preguntas

2. \textbf{Listening} - Escuchar activamente

3. \textbf{Acknowledgment} - Reconocer la pregunta

4. \textbf{Brevity} - Ser conciso (30-90 seg.)

5. \textbf{Calmness} - Mantener la calma

6. \textbf{Honesty} - Admitir cuando no sabes

7. \textbf{Respect} - Respetar todas las preguntas

8. \textbf{Bridge} - Puentear a tu mensaje

\vspace{0.2cm}

---

\vspace{0.2cm}

\section{EJEMPLOS DE RESPUESTAS}

\vspace{0.2cm}

\subsection{Ejemplo 1: Pregunta crítica}

\vspace{0.2cm}

\textbf{P:} "¿No es su propuesta utópica e inviable?"

\vspace{0.2cm}

\textbf{R:} "Entiendo su escepticismo. Es una pregunta legítima. Permítanme presentar la evidencia de casos similares que han tenido éxito. Por ejemplo, en [lugar/caso], implementaron algo similar y lograron [resultado]. Basándonos en esa evidencia, creo que es viable con la planificación adecuada."

\vspace{0.2cm}

\subsection{Ejemplo 2: Pregunta compleja}

\vspace{0.2cm}

\textbf{P:} "¿Cómo afectará su propuesta a los grupos X, Y y Z, considerando las restricciones A y B, y qué pasará si ocurre C?"

\vspace{0.2cm}

\textbf{R:} "Esa es excelente pregunta con múltiples partes. Permítanme abordarlas una por una. Primero, respecto al grupo X... Segundo, sobre el grupo Y... Y tercero, en cuanto a las restricciones..."

\vspace{0.2cm}

\subsection{Ejemplo 3: Pregunta sin respuesta conocida}

\vspace{0.2cm}

\textbf{P:} "¿Cuál es el costo exacto por usuario en 10 años?"

\vspace{0.2cm}

\textbf{R:} "Esa es una excelente pregunta que no puedo responder con precisión ahora. Los modelos proyectan un rango de X-Y, pero dependerá de múltiples factores. Puedo proporcionarles un análisis más detallado después de la presentación."

\vspace{0.2cm}

\subsection{Ejemplo 4: Pregunta hostil}

\vspace{0.2cm}

\textbf{P:} "¿Por qué debería alguien creerle después de sus fracasos anteriores?"

\vspace{0.2cm}

\textbf{R:} [Calmamente] "Veo que hay preocupaciones sobre el historial. Respecto a los proyectos anteriores, [breve contexto]. Lo importante es qué aprendimos y cómo aplicamos esas lecciones al proyecto actual, que tiene [diferencias clave]..."

\vspace{0.2cm}

---

\vspace{0.2cm}

\section{EJERCICIOS PRÁCTICOS}

\vspace{0.2cm}

\subsection{Ejercicio 1: Anticipación}

Escribe 5 preguntas probables que recibirías en tu próxima presentación.

\vspace{0.2cm}

\subsection{Ejercicio 2: Role-play}

En parejas, uno hace una pregunta desafiante y el otro responde usando técnicas de puente o reencuadre.

\vspace{0.2cm}

\subsection{Ejercicio 3: Grabación y análisis}

Graba tu respuesta a una pregunta difícil y luego analiza:

\begin{itemize}

  \item ¿Fui conciso?

  \item ¿Respondí directamente?

  \item ¿Mantuve la calma?

  \item ¿Usé técnicas de puente o reencuadre?

\end{itemize}

\vspace{0.2cm}

---

\vspace{0.2cm}

\section{RÚBRICA DE EVALUACIÓN C1}

\vspace{0.2cm}

\subsection{Excelente (A)}

\begin{itemize}

  \item Responde directamente y concisamente

  \item Maneja preguntas difíciles con calma

  \item Usa técnicas de puente y reencuadre efectivamente

  \item Muestra preparación y conocimiento

  \item Respeta todas las preguntas, incluso las hostiles

  \item Refuerza mensajes clave en respuestas

\end{itemize}

\vspace{0.2cm}

\subsection{Bueno (B)}

\begin{itemize}

  \item Responde adecuadamente la mayoría de preguntas

  \item Maneja bien preguntas moderadas

  \item Algún uso de técnicas de gestión

  \item Muestra preparación básica

  \item Respeta la mayoría de preguntas

  \item Respuestas generalmente claras

\end{itemize}

\vspace{0.2cm}

\subsection{Suficiente (C)}

\begin{itemize}

  \item Responde preguntas básicas

  \item Dificultad con preguntas difíciles

  \item Uso limitado de técnicas

  \item Preparación mínima

  \item Algún desafío con respeto

  \item Respuestas a veces vagas o largas

\end{itemize}

\vspace{0.2cm}

---

\vspace{0.2cm}

\textbf{Sesión 24: La Conferencia (II)}

\textbf{Nivel C1 - Español Oral Avanzado}

\vspace{0.2cm}

\end{document}
